\printedpage{14}
% head: A. Grothendieck : Dérivateurs, Ch. XIII, Édition des Universités de Montpellier II et Paris VII
% foot: Version 25 janvier 2004

\msp{page 24}
\noindent and the inclusions stipulated in 2$^\circ$, $b'$) ensure that we can indeed find $TC$ and $C$, if need be taking them as large as possible, $TC = W^{\univ}$, $C = \Cof_W$ -- in which case, moreover, we will have
\[
TC \;=\; C \cap W^{\univ},
\]
\textit{i.e.} we have a left pre-Q-category of models. But now, on the other hand, the point is to ensure that
\[
TF \;=\; F \cap W_{\univ},
\]
which implies, as we have just seen, (noting $(W^{\univ})_* \subset F$) that
\[
\overline{(W^{\univ})_* \cap W_{\univ}} \;\subset\; \overline{F \cap W_{\univ}} = \overline{TF} = TF \subset W_{\univ}. \tag*{$(*)$}
\]
Conversely, if $(*)$ is satisfied, let us take
\[
F = (W^{\univ})_*
\]
as small as possible, \textit{i.e.} $C = W^{\univ}$ \texttt{[rather} $TC = W^{\mathtt{univ}}$\texttt{]} as large as possible. We must therefore take
\[
TF \;=\; F \cap W_{\univ} \;=\; (W^{\univ})_* \cap W_{\univ},
\]
and this set is indeed right closed by hypothesis.

\textbf{Corollary.} \textit{Suppose condition} b$'$) \textit{of the theorem-scholie is satisfied. The following conditions} a), b) \textit{on $W$ are equivalent (and imply the condition of} 3$^\circ$ \textit{of the theorem-scholie).}

\msp{page 25}
\begin{itemize}[leftmargin=3em]
\item[\cercle{a}] $(W^{\univ})_* \cap \overline{W_{\univ}} = (\Cof_W)_*$.

\item[\cercle{b}] $\widetilde{W^{\univ} \cup (W_{\univ})}^{*} = \Cof_W$, i.e. $\Cof_W$ \textit{is the left closed subset generated by} $W^{\univ}$ \textit{and} $(W_{\univ})^*$\textit{,} % doubt resolved: printed tilde ends before the final star; preserve the edition reading, see corrections.csv (PDF leaf 16).
\end{itemize}

\noindent\textit{and these conditions imply}

\begin{itemize}[leftmargin=3em]
\item[\cercle{c}] \textit{that there exists a structure of a} closed Quillen triple $(W, C, F)$ \textit{including} $W$ \textit{(thus with} $TC = W \cap C$, $TF = W \cap F$\textit{) (which will be closed in Quillen's sense if and only if} $W$ \texttt{[is]} \textit{stable under direct factors -- for $C$ and $F$ are already spontaneously stable).}
\end{itemize}

\noindent Indeed, it suffices then to take
\[
TC \;=\; W^{\univ}, \qquad\quad C \;=\; \Cof_W .
\]
But I doubt that the existence of a closed Quillen triple $(W, C, F)$ suffices to imply conditions \textit{a}), \textit{b}) above, or even the weaker condition
\[
(W^{\univ})_* \cap W_{\univ} \;=\; (\Cof_W)_*
\]
which means \textit{exactly} that $C = \Cof_W$, $TC = W^{\univ}$, and their transposes $TF = (\Cof_W)_*$, $F = (W^{\univ})_*$ form a closed Quillen category.

\clearpage
\printedpage{15}
% head: A. Grothendieck : Dérivateurs, Ch. XIII, Édition des Universités de Montpellier II et Paris VII
% foot: Version 25 janvier 2004

\msp{page 26}
\noindent Conditions on $(\mathcal{M}, W)$ with
\begin{itemize}[leftmargin=3em]
\item[\textit{a})] $\mathcal{M}$ stable under finite $\varprojlim$ and (small) $\varinjlim$,
\item[\textit{b})] $W$ stable under filtering $\varinjlim$,
\item[\textit{c})] $\mathcal{M}$ accessible.
\end{itemize}

\[
\mathrm{I}\ \left\{
\begin{array}{@{}l@{\quad}p{11cm}@{}}
\cercle{a} & existence of a weak Quillen system $(W, C, TC, F, TF)$ (closed if one insists) \\[1ex]
\Updownarrow & \\[1ex]
\cercle{b} & existence of a left exact weak Quillen system ($TC = C \cap W$). \\[1ex]
\Updownarrow & \\[1ex]
\cercle{c} & $\underbrace{(\Cof_W)_* \subset W,\ (W^{\univ})_* \subset \Fib_W}_{(*)}$ ($\Longleftrightarrow$ $(\Cof_W)_* \subset W_{\univ}$, $(W^{\univ})_* \subset \Fib_W$).
\end{array}\right.
\]

\noindent$\Uparrow$
\[
\mathrm{II}\ \left\{
\begin{array}{@{}l@{\quad}p{11cm}@{}}
\cercle{a} & existence of a right exact weak Quillen system ($TF = F \cap W$) \\[1ex]
\Updownarrow & \\[1ex]
\cercle{d} & We have $(*)$, and moreover $(W^{\univ})_* \cap W_{\univ}$ right closed. \\[1ex]
\Updownarrow & \\[1ex]
\cercle{b} & existence of a closed Quillen system (closed model category), $TC = C \cap W$, $TF = F \cap W$, $F = (TC)_*$, $TF = C_*$, $TC = F^*$, $C = (TF)^*$. \\[1ex]
\Updownarrow & \\[1ex]
\cercle{c} & Putting $TC = W^{\univ}$, $F = (TC)_*$, $TF = F \cap W_{\univ}$, $C = (TF)_*$, we find together with $W$ a closed Quillen system. % [sic] $C = (TF)_*$ -> $C = (TF)^*$, cf. the line b) above where $C = (TF)^*$
\end{array}\right.
\]

\noindent$\Uparrow$
\[
\mathrm{III}\ \left\{
\begin{array}{@{}p{11.5cm}@{}}
$C = \Cof_W$ and $TC = W^{\univ}$ are part (together with $W$) of a closed Quillen system (thus $F = (W^{\univ})_*$, $TF = (\Cof_W)_*$, $W = TF \circ TC$). \\[1ex]
$\Updownarrow$ \\[1ex]
\texttt{[We have]} $(*)$, and moreover $(W^{\univ})_* \cap W_{\univ} = (\Cof_W)_*$.
\end{array}\right.
\]

\noindent$\Uparrow$
\[
\mathrm{IV}\ \left\{
\begin{array}{@{}p{11.5cm}@{}}
We have $(*)$ and $\Cof_W = \widetilde{(W_{\univ})^* \cup W^{\univ}}$\ \ . \\[1ex]
$\Updownarrow$ \\[1ex]
We have $(*)$ and $(W^{\univ})_* \cap \overline{W_{\univ}} = (\Cof_W)_*$.
\end{array}\right.
\]

\clearpage
\printedpage{16}
% head: A. Grothendieck : Dérivateurs, Ch. XIII, Édition des Universités de Montpellier II et Paris VII
% foot: Version 25 janvier 2004

\noindent Case $\mathcal{M} = \mathrm{Cat}$, $W = \mathbf{W}_\infty$. We know that II is satisfied, by \textsc{Thomason}. Thus we have
\[
(\Cof_W)_* \subset W_{\univ}, \qquad\quad (W^{\univ})_* \subset \Fib_W, \qquad\quad (W^{\univ})_* \cap W_{\univ} \ \text{ right closed.}
\]
I have been able to prove the first two inclusions directly (\textit{i.e.} I), but not that $(W^{\univ})_* \cap W_{\univ}$ is right closed. (I should at least try!) (\footnote[18]{It would suffice to know that $\overline{W_{\univ}} \cap \Fib_W = W_{\univ}$, which would imply that III and IV are equivalent (\footnotemark[19]).}\footnotetext[19]{But this is false, e.g. if $X$ \texttt{[is a]} nonempty discrete category, $X \longrightarrow e$ is in $\overline{W_{\univ}} \cap \Fib_W$.}). We do not have IV, \textit{i.e.} we have
\[
(W^{\univ})_* \cap \overline{W_{\univ}} \;\neq\; (\Cof_W)_*
\]
for
\[
(W^{\univ})_* \cap \overline{W_{\univ}} \not\subset W_{\univ}
\]
(contains the $X \longrightarrow e$ with $X$ discrete), so we have
\[
(W^{\univ})_* \cap \overline{W_{\univ}} \;\;\underset{\neq}{\supset}\;\; \underbrace{(W^{\univ})_* \cap W_{\univ}}_{\text{is \textit{not} the orthogonal of } W^{\univ} \cup (W_{\univ})^*} \quad . % doubt resolved: printed strict-superset sign confirmed; see corrections.csv (PDF leaf 18).
\]

\msp{page 27}

\section*{3\quad Quillen category (second version)}

We are still given a pair
\[
(\mathcal{M}, W), \qquad\quad W \subset \mathrm{Fl}(\mathcal{M}), \qquad\quad W \text{ weakly saturated,}
\]
$\mathcal{M}$ satisfies M 0 (stability under finite $\varprojlim$ and finite $\varinjlim$).

\noindent\textbf{Definition 1 :} A \textit{weak Quillen structure} associated to the category of models $(\mathcal{M}, W)$ is the datum of four subsets $(C, TC, F, TF)$ of $\mathrm{Fl}(\mathcal{M})$, satisfying the following conditions.

\begin{itemize}[leftmargin=3em,labelsep=0.5em,align=left]
\item[\textbf{M 1}] $C \longleftrightarrow TF$, $TC \longleftrightarrow F$ (\textit{i.e.} $TF \subset C_*$, $F \subset (TC)_*$).
\item[\textbf{M 2}] Quillen factorization property for $(C, TF)$ and $(TC, F)$.
\item[\textbf{M 3}] Stability of $C$ under composition and under cobase change, of $F$ under composition and base change -- $C$ and $F$ contain the isomorphisms.
\item[\textbf{M 4}] Similar stability for $TC$ and $TF$.
\item[$\mathbf{M}'$] $TC \subset C \cap W$, $TF \subset F \cap W$.
\item[$\mathbf{M}''$] $C \subset \Cof_W$, $F \subset \Fib_W$.
\end{itemize}

\noindent We say that the Quillen structure is \textit{exact} if it satisfies the additional axiom

\clearpage
\printedpage{17}
% head: A. Grothendieck : Dérivateurs, Ch. XIII, Édition des Universités de Montpellier II et Paris VII
% foot: Version 25 janvier 2004

\noindent\textbf{M 2 bis}\qquad $W = TF \circ TC$.

\noindent We say that it is \textit{left exact} if it satisfies the additional axiom

\noindent\textbf{M$'_0$ (g)}\qquad $TC = C \cap W$ (instead of only $TC \subset C \cap W$),

\noindent\textit{right exact} if it satisfies

\noindent\textbf{M$'_0$ (d)}\qquad $TF = F \cap W$ (instead of only $TF \subset F \cap W$) (\footnote[20]{\textbf{NB} If $C$, $F$, $W$ \texttt{[are]} stable under direct factors (thus also $C \cap W$ and $F \cap W$), then left exactness $TC = C \cap W$ \textit{is equivalent} to right exactness $TF = F \cap W$, and implies that $C$, $TC$, $F$, $TF$ are \textit{closed} (on the left for $C$, $TC$, on the right for $F$, $TF$).}).

\msp{page 28}
\noindent\textbf{Remark.}
\begin{itemize}[leftmargin=3em]
\item[a)] A left exact or right exact Quillen structure is exact.
\item[b)] If we have a datum $(C, TC, F, TF)$ satisfying M 1, M 2, M$'$, $M''$, when we saturate $C$, $TC$ on the left and $F$, $TF$ on the right, for the stabilities specified in M 2, M 3, we find a Quillen structure. (All the axioms are obvious, except M$'$, for which I refer to p. 9, 10, in particular prop. 4.) If we had M 2 bis, this will still hold (a fortiori) after saturation, \textit{i.e.} we have an exact Quillen structure.
\end{itemize}

\noindent\textbf{Definition 1 bis.} The weak Quillen structure is called a \textit{Quillen structure} (without qualification) if it is left and right exact, or (as one says) ``biexact'', then $(C, F)$ already determines $(C, TC, F, TF)$. (That corresponds exactly to the ``model categories'' of \textsc{Quillen} \texttt{[proper]}.) We say that the weak Quillen structure is \textit{closed} if we have % [sic] determine -> détermine, missing accent
\[
C \clarrow TF\ , \qquad\quad TC \clarrow F\ ,
\]
\textit{i.e.} $TF = C_*$, $C = (TF)^*$ and $F = (TC)_*$, $TC = F^*$ (\footnote[21]{This also means that $C$, $TF$, $TC$, $F$ are stable under direct factors. We say that the weak Quillen structure is \textit{left closed} (resp. \textit{right closed}) if $C$, $TC$ (resp. $F$, $TF$) are left closed (resp. right closed).}). A weak Quillen structure that is both biexact (\textit{i.e.} not weak at all!) and closed is called a \textit{perfect} Quillen structure. (That corresponds exactly to the ``closed model category'' \texttt{[proper]} of \textsc{Quillen}.)

\msp{page 29}
\noindent\textbf{Remark.} For a weak Quillen structure to be perfect, it suffices that it be biexact and that we have
\[
C = (TF)^*, \qquad\quad F = (TC)_*
\]
(which already implies that $TC = F^*$, $TF = C_*$, \textit{i.e.} that $TC$ is left closed and $TF$ right closed). This is proved in Quillen. It also implies that $W$ is strongly saturated -- it's the land of plenty!

\clearpage
\printedpage{18}
% head: A. Grothendieck : Dérivateurs, Ch. XIII, Édition des Universités de Montpellier II et Paris VII
% foot: Version 25 janvier 2004

\noindent\textbf{Scholie :} The pair $TC \longleftrightarrow F$, satisfying the Quillen factorization condition, and satisfying $TC \subset W$, $F \subset \Fib_W$, already allows one to develop a satisfactory theory of $W$-homotopic fibers in $\mathcal{M}$ -- but without an exactness property for $W$-cartesian squares. Dually, the datum of $C \longleftrightarrow TF$, with factorization, and $TF \subset W$, $C \subset \mathrm{Cofib}_W$, gives a good theory of $W$-homotopic cofibers, but without an exactness property for $W$-cocartesian squares. The weak Quillen structure brings together these two data dual to one another, with in addition the additional link between these two data, by $TC \subset C$, $TF \subset F$ (stipulated in $M'$). This datum already ensures that % [sic] donnés -> données, feminine noun as in "ces deux données duales"
\[
\underbrace{\mathcal{M}_{\mathrm{cf}}}_{\substack{\text{full subcategory}\\ \text{of $\mathcal{M}$ consisting of}\\ \text{objects fibrant and}\\ \text{cofibrant at once}}} W_{\mathrm{cf}}^{-1} \;\xrightarrow{\ \sim\ }\; \mathcal{M}W^{-1} = \mathrm{Ho}_W
\]
\msp{page 30}
is an equivalence, but not that the localization functor
\[
\mathcal{M}_{\mathrm{cf}} \longrightarrow \mathcal{M}_{\mathrm{cf}}W_{\mathrm{cf}}^{-1}
\]
is a simple passage to the quotient on arrows, \textit{i.e.} that if $X$, $Y$ are in $\mathcal{M}_{\mathrm{cf}}$, then
\[
\mathrm{Hom}_{\mathcal{M}}(X, Y) \;\longrightarrow\; \underbrace{\mathrm{Hom}_W(X, Y)}_{\overset{\text{def}}{=}\,\mathrm{Hom}_{\mathrm{Ho}_W}(X, Y)} \tag*{$(*)$}
\]
is surjective. This is ensured by the exactness axiom M 2 bis,
\[
W \;=\; TF \circ TC.
\]
With this axiom alone, we have no nice characterization of the equivalence relation in $\mathrm{Hom}_{\mathcal{M}}(X, Y)$ defined by the arrow $(*)$ (for $X$, $Y$ in $\mathcal{M}_{\mathrm{cf}}$).

If we have left exactness,
\[
TC = C \cap W \qquad\quad (\Longleftrightarrow\ TC = C \cap W^{\univ}),
\]
then this equivalence relation is the one generated by the relation of \textit{right Q-homotopy} (\footnote[22]{It would seem that the relation of right Q-homotopy is an equivalence relation, as soon as $Y$ \texttt{[is]} in $\mathcal{M}_{\mathrm{f}}$ and without assuming the weak Quillen structure to be right or left exact.}, \footnote[23]{\textbf{NB} According to \textsc{Baues}, it would seem that on $\mathcal{M}_{\mathrm{cf}}$, this relation of right homotopy is equivalent to the relation of left homotopy, and that the latter does not depend on the cylinder chosen \ldots}) (implied by left Q-homotopy), and the quotient category thus defined
\[
\pi_{\mathrm{d}}\mathcal{M}_{\mathrm{cf}} \;=\; \pi\mathcal{M}_{\mathrm{cf}} \;\simeq\; \mathrm{Ho}_W,
\]
maps into $\pi_{\mathrm{d}}\mathcal{M}_{\mathrm{c}}$,
\[
\pi\mathcal{M}_{\mathrm{cf}} \longrightarrow \pi_{\mathrm{d}}\mathcal{M}_{\mathrm{c}},
\]
and this functor can be interpreted (unless I am mistaken)